\documentclass[10pt,a4paper]{amsart}
\usepackage[margin=19mm]{geometry}
\usepackage{amsmath,amssymb,mathtools}
\usepackage[scr=rsfs,cal=euler]{mathalfa}
\usepackage{xfrac}
\usepackage[unicode,hidelinks,bookmarks=false]{hyperref}
\newcommand{\ie}{i.e.\ }
\newcommand{\cf}{cf.\ }
\newcommand{\resp}{respectively\ }
\newcommand{\Subsection}{\subsection}
\providecommand{\nobreakdash}{\nobreak\hbox{-}}
\setlength{\emergencystretch}{2em}
\multlinegap0pt
\usepackage{mdframed}
\usepackage{mdframed}
\usepackage{mdframed}
\usepackage{mdframed}
\usepackage{mdframed}
\usepackage{enumerate}
\usepackage[shortlabels]{enumitem}
\usepackage{bm}
\usepackage{dsfont}
\usepackage{amssymb}
\usepackage{mathtools}
\usepackage{xcolor}
\usepackage{cancel}
\usepackage{xfrac}
\usepackage{siunitx}
\usepackage{float}
\usepackage{tikz}
\usepackage{stackengine}
\newtheorem{assumption}{Assumption}
\newtheorem{theorem}{Theorem}
\def\P{{\mathrm P}}
\title{\LARGE Nearest-Neighbor Radii under Dependent Sampling}
\author{Yuanyuan Gao, Yilong Hou, Zhexiao Lin}
\date{Source: arXiv:2605.14343v1 (CC BY 4.0)}
\begin{document}
\maketitle

\noindent\emph{Source and licence.} \LARGE Nearest-Neighbor Radii under Dependent Sampling. Version arXiv:2605.14343v1 (CC BY 4.0), distributed by arXiv under the Creative Commons Attribution 4.0 International licence, http://creativecommons.org/licenses/by/4.0/, as recorded in the arXiv licence metadata for this exact version and shown on the abstract page https://arxiv.org/abs/2605.14343v1. Full licence notice: \texttt{licenses/arxiv-2605.14343-CC-BY-4.0.txt}.\par\smallskip

\noindent\textit{Editorial scope.} Counted unit: the article's theorem thm:tail (source0.tex lines 501--516). The article's complete original proof of it is delivered with it (source0.tex lines 1088--1163, one \texttt{proof} environment of 76 lines, which the article places away from the statement). No mathematics is written, repaired or re-derived: the statement and the proof are the article's own lines, in the article's own notation, and no retained line is substituted or re-flowed.  Retained uncounted as the source-local context the proof needs: lines 464--466; lines 472--474; lines 485--487.  Nothing was omitted from any retained span, so the package carries no omission record.  No retained line contains a citation, so the wrapper carries no bibliography and no bibliography entry.  No retained line contains a figure environment, an image-inclusion command or drawing code, so no image is delivered and the package declares no asset dependency.  Editorial formatting only, in addition: an AMS \texttt{amsart} preamble that restates the article's own declarations and macros used by the retained lines, namely the environments \texttt{assumption}, \texttt{theorem} on separate counters, together with the standard packages those lines need.  The retained environments print in this document as Assumption 1, Theorem 1 in order of appearance, whereas the article prints the same items with the numbering of its own full document.  The complete native source of the article as distributed in the arXiv e-print for this exact version is bundled unchanged as \texttt{sources/200578/source0.tex}.
\begin{assumption}[Geometric $\alpha$-mixing]\label{ass:geom}
There exists $\gamma>0$ such that $\alpha_m \le e^{-\gamma m}$ for all $m \ge 1$.
\end{assumption}

\begin{assumption}[Local lower mass of order $s$ at $x$]\label{ass:lower}
There exist constants $c_->0$, $r_0>0$, and $s>0$ such that $\mu(B(x,r)) \ge c_- r^s$ for all $0<r\le r_0$.
\end{assumption}

\begin{assumption}[Compact support]\label{ass:compact}
The support of $\mu$ is contained in a compact set of diameter at most $D$.
\end{assumption}

\begin{theorem}[Non-asymptotic tail bound]\label{thm:tail}
Assume Assumptions~\ref{ass:geom}, \ref{ass:lower}, and \ref{ass:compact}. Define
\[
  b_n := \Big\lceil \frac{8}{\gamma}\log n \Big\rceil,
  \qquad
  r_*(n,k) := \Big(\frac{8k}{c_- n}\Big)^{1/s}.
\]
Then there exists a constant $c_0>0$, depending only on $\gamma$, such that for every integer $j \ge 0$ satisfying $2^j r_*(n,k) \le r_0$, we have
\[
  \P\big(R_{n,k}(x) > 2^j r_*(n,k)\big)
  \le
  4\exp \Big(-c_0\,2^{js}\,\frac{k}{\log n}\Big)
  + C n^{-7},
\]
for all sufficiently large $n$.
\end{theorem}

\begin{proof}

Let $$r_* = \Big(\frac{8k}{c_- n}\Big)^{1/s}$$ and choose \(J_0\) maximal such that $2^{J_0}r_* \le r_0$. Since the support of \(\mu\) has diameter at most \(D\), we have $R_{n,k}(x)\le D$ almost surely.

We decompose the moment over
\[
\{R_{n,k}(x)\le r_*\},
\quad
\{2^j r_* < R_{n,k}(x)\le 2^{j+1}r_*\},\ 0\le j\le J_0-1,
\quad
\{R_{n,k}(x)>2^{J_0}r_*\},
\]
and we obtain
\begin{align*}
\mathbb E[R_{n,k}(x)^p]
&\le r_*^p
+ \sum_{j=0}^{J_0-1}(2^{j+1}r_*)^p\,\mathbb P(R_{n,k}(x)>2^j r_*) + D^p\,\mathbb P(R_{n,k}(x)>2^{J_0}r_*).
\end{align*}

We now apply Theorem~\ref{thm:tail}. For every \(0\le j\le J_0\),
\[
\mathbb P(R_{n,k}(x)>2^j r_*)
\le
4\exp \Big(-c_0\,2^{js}\frac{k}{\log n}\Big)+Cn^{-7},
\]
and then
\begin{align*}
\mathbb E[R_{n,k}(x)^p]
&\le r_*^p
+ C r_*^p \sum_{j=0}^{J_0-1}2^{jp}
\exp \Big(-c_0\,2^{js}\frac{k}{\log n}\Big) + C r_*^p n^{-7}\sum_{j=0}^{J_0-1}2^{jp}
+ C D^p n^{-7}.
\end{align*}

We now bound the three terms on the right-hand side. For the first term, since \(k\ge K_0\log n\), the quantity \(k/\log n\) is bounded below by a positive constant. Therefore
\[
\sum_{j=0}^{J_0-1}2^{jp}
\exp \Big(-c_0\,2^{js}\frac{k}{\log n}\Big)
\le
\sum_{j\ge 0}2^{jp}e^{-c_1 2^{js}}
<\infty
\]
for some constant \(c_1>0\). For the second term, since \(2^{J_0}r_*\le r_0<2^{J_0+1}r_*\), we have $2^{J_0}\asymp r_*^{-1}$, where the constants depend only on \(r_0\). Then
\[
r_*^p\sum_{j=0}^{J_0-1}2^{jp}
\le
C r_*^p 2^{J_0 p}
\le C,
\]
and then
\[
C r_*^p n^{-7}\sum_{j=0}^{J_0-1}2^{jp}
\le C n^{-7},
\]
which is negligible. The third term \(C D^p n^{-7}\) is negligible as well. Combining the above estimates yields
\[
\mathbb E[R_{n,k}(x)^p]
\le C r_*^p + C n^{-7}
\le C r_*^p.
\]
Since
\[
r_*^p
=
\Big(\frac{8k}{c_- n}\Big)^{p/s}
\asymp
\Big(\frac{k}{n}\Big)^{p/s},
\]
we obtain
\[
\mathbb E[R_{n,k}(x)^p]
\le
C\Big(\frac{k}{n}\Big)^{p/s}.
\]
This completes the proof.
\end{proof}

\end{document}
